% -------------------------------------------------------------------------------------
\section{Glauber 算符及其与光锥序证明的对照}\label{sec:scet-glauber}

标准的 SCET$_{\rm I}$、SCET$_{\rm II}$ 拉氏量里没有 Glauber 模式。Bauer--Lange--Ovanesyan 找到一类旁观者相互作用的例子，在其中必须为了自洽而引入 Glauber 胶子\cite{Bauer:2010cc}；Rothstein--Stewart 构造了领头幂次的 Glauber 算符\cite{RothsteinStewart2016}。在图形方法中，Glauber 区是 Drell--Yan 因子化最困难的一步（第~9~章\cref{prop:method-glauber-deformation}，附录~O）；本节说明它在有效理论中的形态。

\begin{proposition}[Glauber 算符 \tagA]\label{prop:scet-glauber-operators}
（Rothstein--Stewart\cite{RothsteinStewart2016}。）(a) 存在耦合 $n$、$\bar n$ 与软扇区的领头幂次 Glauber 算符，SCET$_{\rm II}$ 的 Glauber 拉氏量为\cite[式~(5.21)]{RothsteinStewart2016}
\begin{equation}\label{eq:scet-LG}
\mathcal L_G^{\rm II(0)}=\ee^{-\ii x\cdot\mathcal P}\sum_{n,\bar n}\sum_{i,j=q,g}\mathcal O_n^{iB}\frac{1}{\mathcal P_\perp^2}\mathcal O_s^{BC}\frac{1}{\mathcal P_\perp^2}\mathcal O_{\bar n}^{jC}
+\ee^{-\ii x\cdot\mathcal P}\sum_n\sum_{i,j=q,g}\mathcal O_n^{iB}\frac{1}{\mathcal P_\perp^2}\mathcal O_s^{j_nB},
\end{equation}
% VF-FRO-6 (re-read in the RS source: eq. (5.21))
其中 $\mathcal O_n^{qB}=\bar\chi_nT^B\frac{\bar{\sla n}}{2}\chi_n$，$\mathcal O_n^{gB}=\frac\ii2f^{BCD}\mathcal B^C_{n\perp\mu}\frac{\bar n}{2}\cdot(\mathcal P+\mathcal P^\dagger)\mathcal B^{D\mu}_{n\perp}$，$\mathcal O_{\bar n}$ 同理（$n\leftrightarrow\bar n$），$\mathcal P$ 为标签动量算符；不含软胶子时 $\mathcal O_s^{BC}=8\pi\as\,\delta^{BC}\mathcal P_\perp^2$，于是 $n$ 与 $\bar n$ 共线夸克之间的交换是瞬时势 $-8\pi\as/\bm q_{\mathrm T}^2$（与 $n\cdot q$、$\bar n\cdot q$ 无关，即在 $x^\pm$ 上局域）；BPS 场重定义之前 $\mathcal L_G^{\rm I(0)}=\mathcal L_G^{\rm II(0)}$。
(b) 在硬散射中，Glauber 交换“可以引起使因子化失效的修正”；对硬散射算符，软与共线圈图的 Glauber 零箱减除保证结果对软与共线 Wilson 线的方向不敏感。
(c) 对旁观者—旁观者交换，迭代的 Glauber 交换给出一个相位，当横向动量转移 $\Delta p_\perp$ 的积分不受限制时，它在截面中相消。
(d) 对活跃—活跃交换，在软 Wilson 线取适当方向时，Glauber 交换图等于软图的 Glauber 零箱\cite[第~10.1~节]{RothsteinStewart2016}。
（本命题只涉及部分子散射。）
\uses{def:scet-modes, prop:scet-scet2}
\end{proposition}

\begin{proofsketch}
(1) \emph{树图匹配。}把 QCD 中 $t$ 道单胶子交换的前向散射图投影到领头幂次：对 $n$ 共线夸克 $\bar u_n\gamma^\mu u_n=n^\mu\,\bar u_n\frac{\bar{\sla n}}{2}u_n+\order{\lambda}$，对 $\bar n$ 共线夸克同理；Glauber 动量 $q\sim(\lambda^2,\lambda^2,\lambda)$ 使 $q^2=n\cdot q\,\bar n\cdot q-\bm q_{\mathrm T}^2=-\bm q_{\mathrm T}^2[1+\order{\lambda^2}]$，故
$\ii\mathcal M=-\ii\frac{8\pi\as}{\bm q_{\mathrm T}^2}\big[\bar u_n\tfrac{\bar{\sla n}}{2}T^Bu_n\big]\big[\bar u_{\bar n}\tfrac{\sla n}{2}T^Bu_{\bar n}\big]+\order{\lambda}$，这正是\cref{eq:scet-LG}第一项在 $\mathcal O_s^{BC}=8\pi\as\delta^{BC}\mathcal P_\perp^2$ 时的矩阵元（数值上相对偏差在 $\lambda=10^{-1},10^{-2},10^{-3}$ 时为 $3.2\times10^{-3}$、$2.8\times10^{-5}$、$2.8\times10^{-7}$）；含软胶子的部分由带一个与两个软胶子的匹配确定，其单胶子 Feynman 规则是 Lipatov 顶角\cite[第~5.1、6~节]{RothsteinStewart2016}。% w18_s5 (D1, D2)
(2) \emph{幂次计数。}Glauber 算符是领头幂次的，且其形式没有圈修正\cite[第~5.3~节]{RothsteinStewart2016}（此步给出 E-S 型的陈述）。
(3) \emph{Glauber 圈。}Glauber 圈积分须用快度正规子 $|2k^z|^{-\eta}\nu^\eta$；有了它，交叉箱图为零而箱图给出一个相位（\cref{rem:scet-glauber-dictionary}）\cite[第~5.2~节]{RothsteinStewart2016}。
(4) \emph{硬匹配。}在单圈与两圈显式计算中，活跃线之间的 Glauber 交换给出全部低能的 $\ii\pi$ 项，并与软图的 Glauber 零箱对应，由此得 (b)、(d)\cite[第~10~节]{RothsteinStewart2016}。
(5) \emph{旁观者。}对 Glauber 梯形图求和得到相位 $\ee^{\ii\phi(b_\perp)}$，它在 $\Delta p_\perp$ 的积分不受限制时相消，得 (c)\cite[式~(9.1)、(9.18)--(9.19)、(11.12)--(11.13)]{RothsteinStewart2016}；同一机制见\cite[式~(14)]{AybatSterman2009}。完整论证见\cite[第~5、9--11~节]{RothsteinStewart2016}。
\end{proofsketch}

\paragraph{地位与缺口}
算符的构造（(a)）是全阶论证，依据是匹配与幂次计数；(b)--(d) 是在一圈、两圈与梯形图求和中显式得到的结论，其对所有阶的推广被作者认为“非常可信”但没有证明\cite[第~10~节]{RothsteinStewart2016}。Rothstein--Stewart 本人明确指出：旁观者相位的相消“taken alone it is far from a proof of factorization, even in the abelian case”（单独来看远不是因子化的证明，即使在阿贝尔情形）\cite[第~11.1~节]{RothsteinStewart2016}；SCET 中完整的因子化证明需要证明 Glauber 拉氏量对硬散射过程没有贡献，他们说这将“在别处给出”\cite[第~11~节]{RothsteinStewart2016}——据我们截至 2026 年 10 月的 INSPIRE 检索，这一证明尚未发表。因此最弱的一步是“Glauber 拉氏量的效应在单举截面中相消”本身，它对应附录~O 中的\cref{asm:appO-mixing}(c) 与\cref{asm:appO-deformation}；对 $e^+e^-$ 的 $n$–$\bar n$ 产生流，Rothstein--Stewart 在单圈与两圈给出了显式证据\cite[第~10~节]{RothsteinStewart2016}。这一步是相消问题而不是幂次估计；作为幂次计数，(a) 只给出 E-S 型的陈述。在强子对撞中，Gaunt 发现对强子横能量 $E_T$，CSS 的 Glauber 相消论证在两个旁观者 Glauber 胶子的层次失效，并据此提出对 $E_T$ 等对多部分子相互作用敏感的观测量领头幂次因子化可能失效\cite{Gaunt2014}；Rothstein--Stewart 也认为，当测量约束 $\Delta p_\perp$ 时旁观者 Glauber 交换从 $\order{\as^4}$ 起破坏因子化\cite[第~11.1~节]{RothsteinStewart2016}；Schwartz--Yan--Zhu 则对纯 Glauber 梯形图与任意红外安全的单标度观测量给出了全阶相消的论证，否定了上述论断中纯 Glauber 的部分，而对多重微分的观测量不适用\cite{Schwartz:2018obd}。这些结果的范围与 Glauber 有效理论的进一步状态见第~18~章与第~24~章（\cref{prop:front-glauber-bfkl,prop:front-syz}）。升级途径：在有效理论中证明 $\mathcal L_G$ 对单举观测量的贡献相消，即附录~O 的\cref{thm:appO-main}去掉模型限制后的有效理论版本。

\begin{remark}[Glauber SCET 与光锥序证明的对照]\label{rem:scet-glauber-dictionary}
Rothstein--Stewart 的 Glauber 有效理论与附录~O 的光锥序（LCOPT）证明处理的是同一个动量区域，二者的对应如下（附录~O 的\cref{rem:appO-rs-dictionary}给出同一对照的精确版本）：
\begin{center}\small
\renewcommand{\arraystretch}{1.15}
\begin{tabular}{@{}p{0.43\textwidth}p{0.5\textwidth}@{}}
\toprule
Glauber SCET\cite{RothsteinStewart2016} & 附录~O 的光锥序证明 \\
\midrule
瞬时势 $\mathcal O_n\frac{1}{\mathcal P_\perp^2}\mathcal O_s\frac{1}{\mathcal P_\perp^2}\mathcal O_{\bar n}$ & 等光锥时刻插入与 Glauber 算符 $V_G$（\cref{lem:appO-equal-time,lem:appO-glauber-operator}） \\
Glauber 的“塌缩”到硬顶点 & 与 $V_H$ 重合的接触项（\cref{rem:appO-contact}） \\
旁观者相位在 $\Delta p_\perp$ 积分后相消 & 最大时刻配对与族求和、窗口与整函数载子（\cref{def:appO-windows,lem:appO-carrier}） \\
软与 Glauber 的重叠（Glauber 零箱） & 减除法中的零箱（\cref{def:scet-zero-bin}；附录~N） \\
Glauber 与软模式的混合 & \cref{asm:appO-mixing} \\
快度正规子 $|2k^z|^{-\eta}$：交叉箱图为零 & 共同时序，$p_{\rm opp}=0$（\cref{rem:dy-common-ordering}；\cref{lem:appO-glauber-operator}(c)） \\
\bottomrule
\end{tabular}
\end{center}
最后一行的机制可以直接看到。设 $n$ 共线线与 $\bar n$ 共线线交换两个 Glauber 胶子，圈动量 $k$；在领头幂次两条线的分母分别只含 $n\cdot k=k^0-k^z$ 与 $\bar n\cdot k=k^0+k^z$。箱图中两个分母为 $-\bar n\cdot p\,(n\cdot k)-\Delta_1+\ii0$ 与 $n\cdot p'\,(\bar n\cdot k)-\Delta_2+\ii0$，它们在复 $k^0$ 平面的零点分居实轴两侧；交叉箱图中第二个分母变为 $-n\cdot p'\,(\bar n\cdot k)-\Delta_2+\ii0$，两个零点位于同侧。没有正规子时 $\int\dd k^0\dd k^z$ 是病态的；正规子 $|2k^z|^{-\eta}$ 只依赖于 $k^z$，因而可以在固定 $k^z$ 下先作 $k^0$ 积分：交叉箱图对任意 $\ii0$ 都精确为零，箱图等于 $-\frac{2\pi\ii}{\bar n\cdot p\,n\cdot p'}\,\frac{1}{2k^z-\Delta_1/\bar n\cdot p-\Delta_2/n\cdot p'+\ii0}$，其 $k^z$ 积分的主值部分 $\propto\tan(\pi\eta/2)\to0$，剩下纯虚的 $\delta$ 函数部分——一个相位\cite[第~5.2~节]{RothsteinStewart2016}。% w18_s5 (E1-E3)
在附录~O 的语言中，这正是全部 $s$ 道权重落在非交叉次序上的“共同时序”，它是精确图给出的物理结果，而独立的对称截断给出 $p_{\rm opp}=\frac12$ 的截断假象（\cref{rem:dy-common-ordering}）。Rothstein--Stewart 也强调，这一切单独来看“远不是因子化的证明”。
\end{remark}
